% !TeX root = ../../Book2.tex
\section{延拓算子}
\label{b2:sec:2303}

延拓算子把曲面数据变成环境空间的振荡函数。对偶性把无法直接解释的曲面取值，转化为一个处处有意义的积分，并为 \(TT^*\) 方法提供起点。

\subsection{对偶限制}
\label{b2:subsec:230301}

\begin{proposition}\label{b2:prop:23-duality}
对于 \(1\le p,r<\infty\)，估计 \(R:L^p\to L^r(S)\) 等价于
\(E:L^{r'}(S)\to L^{p'}(\mathbb R^d)\)，两者最佳常数相同。
\end{proposition}
\begin{proof}
对 Schwartz 函数与曲面光滑函数，Fubini 给出
\[
 \int_S\widehat f(\xi)\overline{g(\xi)}\,\mathrm d\sigma
 =\int_{\mathbb R^d} f(x)\overline{Eg(x)}\,\mathrm dx.
\]
一侧估计配以 Hölder 不等式，再对另一侧单位球取上确界，得到反向算子的估计。有限指数时稠密性完成延拓；出现无穷对偶指数时使用通常的 \(L^1\) 对偶 \(L^\infty\) 即可。最佳常数由两次同样操作相等。
\end{proof}
由于曲面有限测度，\(Eg\) 对 \(g\in L^1(S)\) 绝对收敛且有界连续。若 \(g\) 更粗，先在已建立的范数估计下通过稠密极限定义输出；不能把所有延拓数据都假定光滑。

\subsection{Fourier 衰减}
\label{b2:subsec:230302}

光滑曲面密度的 Fourier 衰减由\chapref{b2:ch:21}的驻相公式给出。在非零曲率紧片上，
\(|\widehat{\mathrm d\mu}(x)|\lesssim(1+|x|)^{-(d-1)/2}\)。
若 \(g\) 光滑，则同样结论成立，但常数含有 \(g\) 的导数。\(L^2(S)\) 单位球不控制这些导数，因此不能把这个逐输入衰减直接当作 \(L^2\to L^q\) 算子界。

\ProseParagraphBreak
平移环境变量对应曲面调制：
\[
 E(e^{-2\pi\mathrm ia\cdot\xi}g)(x)=Eg(x-a).
\]
于是固定 \(L^2\) 范数的输入可使输出峰值移到任意远处，排除一个仅由 \(\|g\|_2\) 控制、固定以原点为中心的逐点衰减界。算子估计必须容许这种平移对称性。

\subsection{\texorpdfstring{\(L^p\to L^q\)}{Lp→Lq} 估计}
\label{b2:subsec:230303}

显然 \(\|Eg\|_\infty\le\|g\|_1\)。若已经得到
\(\|Eg\|_{q_0}\le C\|g\|_2\)，曲面有限测度还给出
\(\|Eg\|_\infty\le |S|^{1/2}\|g\|_2\)。对输出两个范数插值，可得到每个 \(q\ge q_0\) 上的 \(L^2\to L^q\) 界。
对输入指数也可与 \(L^1\to L^\infty\) 插值，产生一个三角范围；这些是已证估计的推论，不能据此覆盖全部必要条件范围。

\ProseParagraphBreak
在球面上，曲面函数 \(g=1\) 的尾部检验已经说明输出不属于任意小的 \(q\)。紧曲面只在输入端是有限测度空间，输出端仍然在整个 \(\mathbb R^d\)。将两端都按有限测度嵌入处理会得到错误的范围。

% 扩写定位：b2:subsec:230303
一旦建立 \(L^2(S)\to L^{q_*}\) 的延拓估计，基本插值就给出一片完整的指数区域。另一方面，曲面测度有限，
\(\|Eg\|_\infty\leq\|g\|_1\)。对 \(0\leq\theta\leq1\)，在这两个端点间使用 Riesz--Thorin，得到
\[
 \frac1s=1-\frac\theta2,\qquad
 \frac1q=\frac\theta{q_*},\qquad
 \|Eg\|_q\leq C_\theta\|g\|_s.
\]
故对 \(1\leq s\leq2\)，证明了直线
\(1/q=\frac{d-1}{d+1}(1-1/s)\) 上的估计。再与同一输入指数的 \(L^\infty\) 界插值，可增大输出 \(q\)。对 \(s>2\)，有限曲面测度使 \(L^s(S)\subset L^2(S)\)，于是仍得到所有 \(q\geq q_*\) 的估计。这后一范围来自输入空间包含，不是把刚才的直线延长到 \(\theta>1\)。

\ProseParagraphBreak

对偶地，从 \(L^{p_*}(\mathbb R^d)\to L^2(S)\) 与 \(L^1(\mathbb R^d)\to L^\infty(S)\) 得到
\[
 \frac1p=1-\frac{\theta}{q_*},\qquad \frac1r=\frac\theta2.
\]
这使一个端点定理同时控制较小输入指数和较大曲面输出指数。曲面测度有限，可以把输出指数降低；环境空间测度无穷，却不能据此任意提高输入 \(p\)。这两个方向在实际引用时很容易混淆，必须先区分哪一侧是有限测度空间。

\ProseParagraphBreak

上述区域仍然只是一个已证明的充分区域。Knapp 管的必要条件对更一般 \(s,q\) 仍适用，径向尾部则给出另外的严格限制；两者之间是否还有空缺，需要新的估计。把插值区域画到必要条件所在边界并不会填补该空缺。指数图的作用是帮助核对已知界，而不是替代其证明。

\subsection{\texorpdfstring{\(TT^*\)}{TT*} 表述}
\label{b2:subsec:230304}

在测试函数上，
\[
 ERf=f*\check\sigma,\qquad
 \|Rf\|_2^2=\langle f*\check\sigma,f\rangle.
\]
因此若卷积算子 \(T=f*\check\sigma\) 有界 \(L^p\to L^{p'}\)，就有
\(\|Rf\|_2\le\|T\|^{1/2}\|f\|_p\)。
反过来，限制界与其对偶延拓界复合也给出同一个卷积界。这里 \(TT^*\) 的实质是一个正的二次型，即便卷积核本身不是非负函数。

\ProseParagraphBreak
若只取核绝对值并使用 Hardy–Littlewood–Sobolev，核的衰减阶 \((d-1)/2\) 会给出较差的指数。Stein–Tomas 的改进保留核的振荡和 Fourier 侧的曲面结构。下一节用解析族把这两类信息放到不同边界线上，再插值合并。

球面延拓还有一个简单的偏微分方程解释。对 \(g\in L^2(S^{d-1})\)，有限曲面测度使 \(g\in L^1(S)\)，因此可以逐阶对 \(x\) 求导，并有全部导数的统一上界。直接计算得到
\[
 (-\Delta-4\pi^2)Eg=0.
\]
这是一族光滑的 Helmholtz 解，频率恰好集中于单位能量面。光滑性来自频率支集紧致，而衰减及全空间可积性则需要曲率和输入性质；两者不可混同。

\ProseParagraphBreak

任何非零的这种解都不属于全空间 \(L^2\)。若 \(Eg\in L^2\)，其 Fourier 变换由 Plancherel 为一个 \(L^2\) 密度；另一方面，温和分布 Fourier 计算给出 \(g\,\mathrm d\sigma\)，支于 Lebesgue 零测的球面。一个绝对连续平方密度若只支于这个零测集只能为零，故 \(g=0\)。如果 \(Eg\in L^q\)、\(q<2\)，又因它有界，插值也会给出 \(L^2\)，同样不可能。限制估计的对偶输出因而自然位于大于二的指数，而不是一个通常的平方能量空间。

\ProseParagraphBreak

这也解释了 \(TT^*\) 中的平方空间究竟在哪里：平方范数落在曲面数据上，环境输出使用较大的 \(q\)。将两者都误写成全空间 \(L^2\)，会把曲面测度与 Lebesgue 测度之间的区别抹掉。用于散射和谱问题时，球面延拓是一种固定能量的广义波，而不是一个具有有限全空间平方质量的束缚态。
